% ================================================================
% CONCLUSION
% ================================================================

% TODO: Summarize main contributions:
%
% 1. Framework connecting targeted CPP to aggregate emissions through
%    production networks.
%
% 2. The I-O structure matters: it can amplify or dampen the direct
%    abatement effects of targeted CPP, governed by the covariance
%    between firms' policy exposure and network-adjusted emission
%    intensity.
%
% 3. Empirical evidence from Belgium suggesting the I-O structure
%    amplifies the effects of the current EU ETS targeting.
%
% 4. Counterfactual exercises showing that targeting firms with
%    higher network centrality can improve emissions outcomes.
%
% TODO: Discuss limitations and directions for future work:
%   - Open economy / trade (carbon leakage across borders)
%   - Dynamic considerations (entry/exit, technology adoption)
%   - Political economy of targeting
%   - General equilibrium effects on output with heterogeneous agents
%   - Richer substitution patterns (non-CES)
